\documentclass[12pt]{article}
\usepackage[margin=1in]{geometry}
\usepackage{hyperref}

\title{\textbf{Problem Set 5}: Web Scraping and APIs}
\author{Jhinuk Banerjee}
\date{\today}

\begin{document}
\maketitle

\section{Web Scraping: Perfume Data from Wikipedia}

\subsection*{What data did I collect?}
I scraped a table for comparing the natural versus the synthetic perfume
ingredients from a Wikipedia article on Perfumes
(\url{https://en.wikipedia.org/wiki/Perfume}).
This table contains five rows comparing natural and synthetic perfumes 
across different dimensions like variance, components, scent uniqueness,
scent complexity, and price. I used the \texttt{rvest} and
\texttt{httr} packages in R to send an HTTP GET request with a
browser user-agent header, parse the returned HTML, and extract
the \texttt{wikitable}-class table into a data frame.

\subsection*{Why does this interest me?}
I really like perfumes and thought it would be interesting to get data on something I like. The natural versus synthetic divide in the
fragrance industry reflects different tensions in luxury markets, like authenticity, relevance for
consumption and gender-targeted marketing. The price
row is interesting from an economics perspective,
as it reflects how the extraction methods and supply chain constraints
shape luxury goods pricing.

\subsection*{Will it be useful?}
This table is quite illustrative, and the scraping skills practiced here can be used to collect larger fragrance datasets. For instance, getting cross-sectional data on perfume characteristics,
gender targeting and pricing could be useful for future research
on gender and consumption.

\subsection*{Resources used}
I used the \texttt{rvest} tidyverse documentation
(\url{https://rvest.tidyverse.org/}) and the \texttt{httr} package
documentation for guidance on setting user-agent headers to
handle HTTP connections reliably.

\section{API Data: World Bank Gender Indicators}

\subsection*{What data did I collect?}
I used the \texttt{wbstats} package in R to query the World Bank
API. I retrieved two indicators for all countries from 2000 to
2023: (1) female labor force participation as a share of the total
labor force (\texttt{SL.TLF.TOTL.FE.ZS}), and (2) the share of
seats held by women in national parliaments
(\texttt{SG.GEN.PARL.ZS}). The resulting dataset contains 5,208
country-year observations.

\subsection*{What did I find?}
The data reveal significant cross-national variation in both
indicators. Female labor force participation ranges from a minimum
of about 6.8\% to a maximum of 54.8\%, with a mean of 41.1\%.
Women's parliamentary representation ranges from 0\% to 63.75\%,
with a mean of 18.8\%, reflecting that women remain
underrepresented in legislatures globally, even though their participation
has grown over the period. 

\subsection*{Why is this useful?}
These indicators are directly relevant to my dissertation research
on economic insecurity and gender in politics. The World Bank
panel provides clean, standardized cross-national data that could
complement my U.S. state legislature elections data. 

\subsection*{Packages used}
I used the \texttt{wbstats} package (v1.1), which provides a clean
R interface to the World Bank's public API. No API key is required.
Documentation is available at
\url{https://CRAN.R-project.org/package=wbstats}.

\end{document}